\chapter{REVISIÓN DE LA LITERATURA}

La calidad de los ambientes interiores (IEQ, por sus siglas en inglés), particularmente en entornos educativos, es un factor crítico que afecta directamente el bienestar, la salud, la seguridad y la comodidad de sus ocupantes \cite{lopez_carreno_2025, teba_2021}. Este capítulo presenta una revisión de la literatura científica y técnica relacionada con la evaluación de la calidad ambiental interior, con un enfoque específico en los sistemas de monitoreo de niveles de CO\(_2\) y otros parámetros relevantes en espacios universitarios, así como la integración de tecnologías emergentes como el Internet de las Cosas (IoT) y LoRaWAN. Se exploran los antecedentes en la medición de la calidad del aire, las normativas nacionales e internacionales, los índices de calidad ambiental, y se examinan soluciones recientes de monitoreo basadas en IoT, con énfasis en el uso de tecnologías de comunicación inalámbrica como LoRa y las plataformas de software para la recolección, visualización y análisis de datos. Este análisis permitirá identificar las brechas existentes y fundamentar el diseño propuesto en esta tesis.

\section{Antecedentes de la Calidad Ambiental Interior y su Monitoreo}

\subsection{Definición e Importancia de la Calidad Ambiental Interior}
La calidad del aire de interiores (CAI), un componente crucial de la IEQ, es determinante para la salud respiratoria y el confort de las personas \cite{sandoval_etal_2020}. Más allá de los agentes patógenos, existen otros factores que impactan el bienestar y que pueden ser medidos y controlados \cite{teba_2021}. La norma UNE 171330:2008 define la Calidad Ambiental en Interiores (CAI) como las condiciones ambientales de los espacios interiores adecuadas para el usuario y la actividad, determinadas por los niveles de contaminación química, microbiológica y los valores de los factores físicos, excluyendo recintos de uso industrial o agrícola \cite{teba_2021}.

Altas concentraciones de contaminantes como el dióxido de carbono (CO\(_2\)) en espacios cerrados pueden generar impactos negativos en la salud y aumentar el riesgo de transmisión de enfermedades respiratorias \cite{contreras_perez_2020, effects_of_low_level_inhalation_exposure_to_carbon_dioxide_in_indoor_environments_a_short_review_on_human_health_and_psychomotor_performance}. El CO\(_2\), proveniente principalmente de la respiración humana en espacios no industriales, es un indicador clave de la ventilación \cite{contreras_perez_2020}. Niveles bajos de humedad también pueden favorecer la transmisión de virus, mientras que la temperatura es crucial para el confort y la salud en espacios concurridos \cite{contreras_perez_2020}. La pandemia de COVID-19 resaltó la importancia de la ventilación adecuada, siendo la circulación de aire exterior el mecanismo más efectivo para reducir la concentración de aerosoles potencialmente infecciosos en espacios como aulas \cite{contreras_perez_2020}. La evaluación de la CAI es, por tanto, fundamental para gestionar riesgos y mejorar la experiencia y salud de los ocupantes \cite{contreras_perez_2020, envira_nodate, higiene_ambiental_nodate}.

\subsection{Evolución Histórica del Monitoreo de la Calidad del Aire}
El monitoreo de la calidad del aire ha progresado significativamente desde mediados del siglo XX. En los años 1940, los primeros dispositivos eran rudimentarios, como tiras de caucho que medían ozono por su agrietamiento, pero requerían mantenimiento constante y tenían precisión limitada \cite{The_Evolution_of_Air_Quality_Monitoring}. Posteriormente, se usaron medidores de CO\(_2\) basados en soluciones químicas, aunque con problemas de fugas y calibración \cite{The_Evolution_of_Air_Quality_Monitoring}. La implementación de redes de estaciones de monitoreo continuo, como la de Los Ángeles en los años 50, permitió la recolección continua de datos y la emisión de alertas, sentando las bases para los sistemas modernos \cite{The_Evolution_of_Air_Quality_Monitoring, indoor_air_quality_guidelines_from_across_the_world_an_appraisal_considering_energy_saving_health_productivity_and_comfort}.

En décadas recientes, la tecnología ha permitido sistemas más compactos y eficientes, utilizando luz ultravioleta, electrónica de estado sólido y sensores avanzados. La irrupción de las redes de sensores inalámbricos (WSN) y el Internet de las Cosas (IoT) ha revolucionado la recolección de datos en tiempo real para contaminantes como CO, NO\(_2\), ozono y CO\(_2\), transmitiendo datos a través de redes interconectadas \cite{Indoor_Air_Quality_Monitoring_Systems_A_Comprehensive_Review_of_Different_IAQM_Systems, The_Evolution_of_Air_Quality_Monitoring}. Estos sistemas, apoyados por el Big Data, gestionan grandes volúmenes de información y facilitan el acceso a resultados mediante portales web y aplicaciones móviles, permitiendo alertas automáticas y una instalación de bajo costo \cite{Indoor_Air_Quality_Monitoring_Systems_A_Comprehensive_Review_of_Different_IAQM_Systems}.

\subsection{Regulaciones y Estándares de Calidad Ambiental Interior}
Las regulaciones sobre IEQ, y específicamente sobre los niveles de CO\(_2\), varían internacionalmente, siendo el CO\(_2\) un indicador clave de la ventilación adecuada. Aunque la Organización Mundial de la Salud (OMS) no ha fijado normas específicas para CO\(_2\) en interiores, muchas directrices nacionales recomiendan umbrales de 1000 ppm para una buena calidad del aire \cite{indoor_air_quality_guidelines_from_across_the_world_an_appraisal_considering_energy_saving_health_productivity_and_comfort, ISIAQ_database}. Niveles entre 1000 y 1500 ppm se consideran aceptables pero moderados, y superiores a 1500 ppm indican calidad deficiente y ventilación insuficiente \cite{indoor_air_quality_guidelines_from_across_the_world_an_appraisal_considering_energy_saving_health_productivity_and_comfort}. En entornos laborales, regulaciones como las de EE. UU. fijan límites de 5000 ppm para CO\(_2\) por toxicidad \cite{indoor_air_quality_guidelines_from_across_the_world_an_appraisal_considering_energy_saving_health_productivity_and_comfort}. ASHRAE Standard 62.1 sugiere que los niveles en edificios públicos no superen los 1000 ppm por encima de los niveles exteriores \cite{indoor_air_quality_guidelines_from_across_the_world_an_appraisal_considering_energy_saving_health_productivity_and_comfort, A_Systematic_Review_of_Air_Quality_Sensors_Guidelines_and_Measurement_Studies_for_Indoor_Air_Quality_Management}.

A nivel europeo, existe legislación sobre ventilación y concentración de CO\(_2\) \cite{teba_2021}. En España, el Reglamento de Instalaciones Térmicas de los Edificios (RITE) establece exigencias de calidad del aire interior e inspecciones anuales en edificios de uso público \cite{teba_2021}. El Código Técnico de Edificación (CTE) también define requisitos y concentraciones máximas de contaminantes \cite{teba_2021}. Estándares como UNE 171330-1:2008 y RESET Air proporcionan directrices para la inspección, control y monitorización continua de la IEQ \cite{teba_2021}. En Perú, la Resolución Ministerial N.º 022-2024-MINSA y su Directiva Administrativa N.º 349-MINSA-DIGIESP-2024 establecen disposiciones para la vigilancia, prevención y control de la salud de los trabajadores con riesgo de exposición a SARS-CoV-2, lo que indirectamente subraya la importancia de la ventilación y calidad del aire en espacios cerrados \cite{directiva_minsa}.

\section{Tecnologías Emergentes para el Monitoreo de la IEQ}

\subsection{Advenimiento del Internet de las Cosas (IoT) en el Monitoreo Ambiental}
El Internet de las Cosas (IoT) se refiere a la interconexión de objetos cotidianos equipados con sensores, software y conectividad, permitiéndoles comunicar datos entre sí y con plataformas centralizadas \cite{contreras_perez_2020, The_Internet_of_Things_A_survey}. Las primeras aplicaciones en monitoreo ambiental surgieron con la integración de sensores inalámbricos, dispositivos portátiles y etiquetas RFID, permitiendo la recolección de datos en tiempo real sin intervención constante \cite{On_the_use_of_IoT_and_Big_Data_Technologies_for_Real_time_Monitoring_and_Data_Processing}. Estas tecnologías posibilitaron el seguimiento continuo de variables como temperatura, ubicación y niveles de contaminación, siendo ampliamente aceptadas en el sector del monitoreo ambiental \cite{mokolo_2022, IAQ_monitoring_systems_based_on_internet_of_things_a_systematic_review}. El IoT ha tenido un impacto significativo en la logística, salud (rastreo de equipos, datos de pacientes) \cite{The_Internet_of_Things_A_survey, the_application_of_internet_of_things_in_healthcare_a_systematic_literature_review_and_classification}, y se ha expandido a la educación y entornos inteligentes, mejorando la seguridad y el consumo energético \cite{On_the_use_of_IoT_and_Big_Data_Technologies_for_Real_time_Monitoring_and_Data_Processing}. Un ejemplo temprano fue el monitoreo de calidad del aire urbano mediante WSN para medir contaminantes y transmitir datos para análisis, optimizando la ventilación y mejorando la salud pública \cite{Developed_urban_air_quality_monitoring_system_based_on_wireless_sensor_networks}.

\subsection{Sistemas de Monitoreo de IEQ Basados en IoT y Tecnologías de Comunicación}
Los sistemas de monitoreo de IEQ basados en IoT permiten obtener datos fiables y en tiempo real sobre diversos contaminantes (CO\(_2\), PM\textsubscript{2.5}, VOCs) mediante sensores de bajo costo, facilitando el monitoreo continuo a largo plazo y la optimización de estrategias de ventilación \cite{Achieving_better_indoor_air_quality_with_IoT_systems_for_future_buildings_Opportunities_and_challenges, indoor_air_quality_assessment_using_a_co2_monitoring_system_based_on_internet_of_things}. Un sistema típico requiere software de procesamiento, estaciones de monitorización con sensores y una interfaz para el usuario \cite{contreras_perez_2020}.

En \cite{Opensource_Based_IoT_Platform_and_LoRa_Communications_with_Edge_Device_Calibration_for_Real_Time_Monitoring_Systems}, se describe un sistema con la plataforma de código abierto BKThings para gestionar dispositivos y visualizar datos en tiempo real de sensores como MQ-135 (calidad del aire) y DHT11 (temperatura y humedad). La transmisión usa LoRa y MQTT a una pasarela Raspberry Pi.
Otro estudio \cite{Achieving_better_indoor_air_quality_with_IoT_systems_for_future_buildings_Opportunities_and_challenges} destaca el uso de Wi-Fi (alta velocidad) y ZigBee (bajo consumo) para la comunicación. Los sensores infrarrojos no dispersivos (NDIR) son comunes para CO\(_2\) por su precisión. Algunas plataformas integran modelos de predicción basados en aprendizaje automático y permiten visualización en tiempo real y descarga de datos.

\subsection{Tecnología LoRa y LoRaWAN para Redes de Monitoreo en Campus Universitarios}
La tecnología LoRa (Long Range) es una modulación de espectro ensanchado que permite comunicación inalámbrica de largo alcance y bajo consumo, ideal para aplicaciones IoT \cite{mokolo_2022}. LoRaWAN es un protocolo de red de área amplia de baja potencia (LPWAN) que utiliza la capa física LoRa y es adecuado para redes extensas como las de un campus universitario, asegurando transmisión segura mediante cifrado AES \cite{Low_Cost_LoRa_Based_IoT_Edge_Device_for_Indoor_Air_Quality_Management_in_Schools}.

En \cite{Low_Cost_LoRa_Based_IoT_Edge_Device_for_Indoor_Air_Quality_Management_in_Schools}, se propone un sistema LoRaWAN para gestionar IAQ en escuelas, midiendo CO\(_2\) (estimado con SGP30), VOCs (SGP30), PM (SPS30), temperatura y humedad (DHT22). Los datos se transmiten a un servidor en la nube y se visualizan en Grafana.
La viabilidad de implementar redes LoRaWAN en campus universitarios se ha estudiado para reducir interferencias, maximizar alcance y disminuir costos \cite{universidad_sevilla_nodate, masache_cabrera_2022}. Un proyecto piloto en un campus buscó desplegar una red LoRa para recoger datos de sensores y transmitirlos de forma segura para crear un campus inteligente y sostenible \cite{alai_secure_nodate}. LoRa es considerada una opción adecuada y costo-efectiva para soluciones de monitoreo ambiental en estos entornos \cite{mokolo_2022}.


\section{Estrategias de mitigación activa del CO$_2$ en ambientes interiores}
\subsection{Ventilación controlada por concentración de CO$_2$}

La concentración de dióxido de carbono en ambientes interiores no solo constituye un indicador del nivel de ventilación, sino que también puede emplearse como variable de control para regular sistemas de renovación de aire. En este enfoque, conocido como ventilación controlada por demanda basada en CO$_2$, el caudal de aire exterior o la capacidad de extracción se ajusta en función de la ocupación real del espacio y de la concentración medida en el ambiente. Este tipo de estrategia ha sido estudiado y aplicado en edificios comerciales, institucionales y educativos, debido a su potencial para mantener condiciones aceptables de calidad del aire interior y, al mismo tiempo, optimizar el consumo energético del sistema de ventilación.

\subsection{Sistemas de mitigación activa mediante ventiladores o extractores}

A diferencia de los sistemas orientados únicamente al monitoreo, los sistemas de mitigación activa buscan reducir la concentración de CO$_2$ mediante la renovación del aire interior. En espacios como aulas, oficinas y otros recintos cerrados, esta reducción puede lograrse a través de ventilación natural, ventilación mecánica o extracción forzada. Diversos estudios han mostrado que la configuración del flujo de aire y la ubicación de extractores o ventiladores influyen directamente en la disminución de los niveles de CO$_2$ y en la efectividad de ventilación del ambiente.

En entornos educativos, la ventilación mecánica descentralizada ha demostrado ser una alternativa útil para mantener concentraciones de CO$_2$ más bajas que las obtenidas únicamente mediante apertura manual de ventanas. Esto resulta especialmente relevante en aulas con alta ocupación o con limitaciones arquitectónicas que dificultan la ventilación natural continua.

\subsection{Soluciones descentralizadas y de bajo costo para ambientes interiores}

Gran parte de las soluciones reportadas en la literatura se implementan sobre infraestructura HVAC existente o requieren sistemas de ventilación mecánica de mayor escala. Sin embargo, en espacios donde no existe infraestructura especializada o donde el presupuesto es limitado, surge la necesidad de desarrollar propuestas descentralizadas de menor costo, compuestas por sensores de CO$_2$, lógica de control y actuadores de ventilación de fácil integración.

En este contexto, el uso de ventiladores controlados por PWM representa una alternativa viable para modular el caudal de aire de acuerdo con la concentración medida en el ambiente. Este enfoque permite plantear estrategias de mitigación activa accesibles y escalables, especialmente en espacios universitarios o académicos donde se busca complementar el monitoreo con acciones correctivas automáticas.

\subsection{Brecha identificada y aporte de la tesis}

A partir de la revisión de la literatura, se observa que muchas propuestas se concentran en el monitoreo de la calidad ambiental interior, la transmisión de datos en tiempo real y la visualización de variables como CO$_2$, temperatura y humedad. No obstante, existe menor desarrollo de soluciones de bajo costo que integren, en un mismo sistema, monitoreo, comunicación IoT y mitigación activa mediante ventiladores controlados automáticamente en función de la concentración de CO$_2$.

En ese sentido, la presente tesis busca complementar el enfoque de monitoreo con un subsistema de actuación basado en ventiladores PWM, orientado a reducir la concentración de CO$_2$ en espacios universitarios de manera automática y económicamente accesible. Esta integración pretende aportar una solución práctica para mejorar la calidad ambiental interior más allá de la simple generación de alertas.



\section{Índices de Calidad Ambiental Interior y Aplicaciones en Entornos Universitarios}

\subsection{Desarrollo y Aplicación de Índices de IEQ}
Para simplificar la interpretación de múltiples parámetros ambientales, se han desarrollado índices de calidad ambiental interior. Un ejemplo relevante es el "Classroom Indoor Air Quality (CIAQ) Risk Index", diseñado para evaluar la capacidad de las aulas para mantener los niveles de CO\(_2\) dentro de límites aceptables \cite{lopez_carreno_2025}. Este índice considera la probabilidad de superar umbrales de CO\(_2\), el número potencial de individuos expuestos y la capacidad del aula para mitigar amenazas. Sirve como herramienta de cumplimiento, ayuda a identificar estrategias de mitigación y facilita la identificación proactiva de riesgos para mejorar la calidad de vida en comunidades educativas \cite{lopez_carreno_2025}. La evaluación de la IEQ y la implementación de un plan de gestión pueden incluir diagnóstico inicial, evaluación de riesgos, diseño de muestreo, acciones correctivas y seguimiento \cite{higiene_ambiental_nodate}.

\subsection{Estudios de Caso y Proyectos en Entornos Educativos y Universitarios}
Diversos estudios han aplicado tecnologías de monitoreo de IEQ en entornos educativos. Un estudio en una universidad evaluó la calidad del aire en aulas midiendo CO\(_2\), humedad y temperatura con sensores conectados a un microcontrolador, enviando datos vía IoT a la plataforma en la nube \texttt{thinger.io} para visualización y análisis \cite{contreras_perez_2020}. El objetivo era identificar aulas con deficiente calidad del aire y proponer mejoras en la ventilación. Se encontraron salones con concentraciones de CO\(_2\) que superaban los límites admisibles, lo que podría causar molestias y aumentar riesgos \cite{contreras_perez_2020}.
Otro estudio evaluó la CAI en aulas de una institución de educación superior midiendo temperatura, humedad relativa, PM\textsubscript{10}, PM\textsubscript{2.5}, CO\(_2\), CO y HCHO, encontrando áreas de mejora, aunque el CO\(_2\) se mantuvo dentro de límites regulatorios en ese caso \cite{sandoval_etal_2020}. Estos trabajos subrayan la necesidad de sistemas de monitoreo continuo para asegurar ambientes saludables en campus universitarios, donde las personas pasan una cantidad significativa de tiempo \cite{its_about_time_a_comparison_of_canadian_and_american_time_activity_patterns}. La arquitectura sostenible también integra la calidad del aire como concepto básico de salud, con estándares como LEED, WELL y BREEAM que pueden ser apoyados por sistemas de monitoreo en tiempo real \cite{envira_nodate}.

% --- FIN DEL CAPÍTULO ---